%=========================================================================
%  ANÁLISIS DISTRIBUCIONAL DEL EXAMEN DE ADMISIÓN UNAM 2026
%  Andamiaje / plantilla — el contenido lo redacta el autor.
%  Compilar en Overleaf con pdfLaTeX.
%=========================================================================
\documentclass[11pt,a4paper]{article}

% ---- Paquetes esenciales ----
\usepackage[utf8]{inputenc}      % acentos y ñ directos en el fuente
\usepackage[T1]{fontenc}         % codificación de salida correcta
\usepackage[spanish]{babel}      % reglas tipográficas en español
\usepackage{amsmath,amssymb}     % matemática: \sum, \log, ecuaciones
\usepackage{graphicx}            % insertar figuras (tus PNG)
\usepackage{booktabs}            % tablas profesionales (\toprule, \midrule)
\usepackage{siunitx}             % alinear números por el punto decimal
\usepackage[hidelinks]{hyperref} % enlaces y referencias clicables
\usepackage{caption}
\usepackage{geometry}
\geometry{margin=2.5cm}

% ---- Metadatos ----
\title{\textbf{TÍTULO — tu pregunta central reformulada aquí}}
\author{Emmanuel}
\date{\today}

\begin{document}
\maketitle

%-------------------------------------------------------------------------
\begin{abstract}
% GUÍA (borra esto al redactar):
% 4-6 oraciones. Un resumen NO es una introducción: aquí va TODO comprimido
% —problema, qué hiciste, qué encontraste (los números clave), qué concluyes.
% Se escribe AL FINAL, cuando ya sabes qué dice el resto. Déjalo vacío por ahora.
\end{abstract}

%-------------------------------------------------------------------------
\section{Introducción}
% GUÍA: aquí van TRES cosas, en este orden:
%   1. El contexto del caso (examen en línea 2026, la polémica, qué pasó).
%   2. Tu PREGUNTA CENTRAL, planteada como hipótesis rivales.
%   3. Por qué importa / qué aporta este análisis (datos abiertos, reproducible).
% NO metas resultados ni matemática todavía. Solo el "qué" y el "por qué".

%-------------------------------------------------------------------------
\section{Marco teórico y fundamentos matemáticos}
% GUÍA: decidiste poner los fundamentos DENTRO del formato canónico, aquí,
% antes de Métodos. Esta sección explica las HERRAMIENTAS, no los datos.

\subsection{Divergencia de Kullback--Leibler}
% GUÍA: define qué mide la KL conceptualmente (divergencia entre una
% distribución P y una referencia Q). Presenta la fórmula. Explica la
% ASIMETRÍA (por qué D(P||Q) != D(Q||P)) y qué dirección elegiste y por qué.
% Menciona la unidad (nats, por el log natural).
%
% Ecuación lista para que la uses (ajusta la notación a tu gusto):
\begin{equation}
    D_{\mathrm{KL}}(P \,\|\, Q) = \sum_{i} P(i) \, \log \frac{P(i)}{Q(i)}
    \label{eq:kl}
\end{equation}

\subsection{Prueba de Kolmogorov--Smirnov}
% GUÍA: define el estadístico D como la máxima distancia vertical entre dos
% CDF empíricas. Aclara qué hipótesis nula prueba. Explica —y esto es clave
% que lo entiendas bien— por qué con n grande el p-valor pierde utilidad y
% por qué reportas D como tamaño de efecto.
\begin{equation}
    D = \sup_{x} \, \lvert F_{1}(x) - F_{2}(x) \rvert
    \label{eq:ks}
\end{equation}

%-------------------------------------------------------------------------
\section{Datos y métodos}
% GUÍA: aquí SÍ hablas de los datos concretos. Subsecciones sugeridas:

\subsection{Fuente y descripción de los datos}
% Origen (datos abiertos DGAE, repositorio comunitario), tamaño (1{,}027{,}716
% folios), rango (2021-2026), las 4 áreas, escala 0-120 aciertos.
% Sé honesto sobre que es scraping de terceros (limitación).

\subsection{Limpieza y criterios de inclusión}
% Tus decisiones EXPLÍCITAS: por qué eliminaste aciertos nulos (no presentaron)
% y ceros; por qué NO filtraste por 'acreditado' (su nulo = "no aplica", no
% dato perdido — el punto que te salvó de sesgar la muestra). Da el n final.

\subsection{Procedimiento de cálculo}
% Cómo discretizaste (bins de 1 acierto), el suavizado épsilon anti-log(0),
% baseline conjunto 2021-2025 como Q vs 2026 como P, y el análisis de
% contraste año-a-año que da escala a la KL. Menciona librerías (scipy).

%-------------------------------------------------------------------------
\section{Resultados}
% GUÍA: SOLO los hallazgos, sin interpretarlos aún (la interpretación va en
% Discusión). Aquí van tus tablas y figuras con su lectura factual.

\subsection{Estadística descriptiva y forma de la distribución}
% Media/mediana por año, y el hallazgo de forma: skewness estable ~+0.8 que
% colapsa a ~0 en 2026, kurtosis que cae a negativa. Aquí va tu figura de
% las 6 distribuciones.

% Ejemplo de inserción de figura (descomenta y ajusta la ruta):
% \begin{figure}[ht]
%     \centering
%     \includegraphics[width=\textwidth]{dist_areas.png}
%     \caption{Distribuciones de aciertos por área: 2026 vs histórico.}
%     \label{fig:dist}
% \end{figure}

\subsection{Divergencia global y contraste histórico}
% KL global 0.339 nats vs 0.002-0.004 año-a-año. KS D=0.293 vs 0.009-0.031.
% Aquí conviene una tabla. Esqueleto listo:
\begin{table}[ht]
    \centering
    \caption{Métricas de divergencia: 2026 vs.\ variación histórica.}
    \begin{tabular}{lccc}
        \toprule
        Métrica & 2026 vs.\ histórico & Rango año-a-año & Factor \\
        \midrule
        KL (nats)      & 0.339 & 0.002--0.004 & $\sim$80--160$\times$ \\
        KS ($D$)       & 0.293 & 0.009--0.031 & $\sim$10--30$\times$ \\
        \bottomrule
    \end{tabular}
    \label{tab:global}
\end{table}

\subsection{Análisis por área}
% Tu tabla de las 4 áreas (KL, KS D, salto de mediana, %>=100). El hallazgo
% de que ambas métricas rankean igual y que el efecto escala con la
% competitividad de las carreras.

%-------------------------------------------------------------------------
\section{Discusión}
% GUÍA: AQUÍ interpretas. Retoma tu pregunta central de la Introducción y
% respóndela con la evidencia. El punto fino que discutimos: la KL mide
% MAGNITUD, no causa; lo que inclina la balanza entre "mejora generacional"
% y "efecto externo" es la FORMA del cambio (segunda población en la cola
% alta, no desplazamiento uniforme). Sé explícito sobre qué puede y qué NO
% puede concluir tu método.

\subsection{Limitaciones}
% Datos de terceros vía scraping; KL/KS no identifican mecanismo ni culpables;
% examen de control presencial no incluido; análisis a nivel área (no
% carrera-plantel); área de Humanidades con n menor.

%-------------------------------------------------------------------------
\section{Conclusión}
% GUÍA: 1-2 párrafos. Cierra respondiendo la pregunta con la fuerza que la
% evidencia permite —ni más ni menos. Sin números nuevos que no hayan
% aparecido antes. Qué queda establecido y qué queda como pregunta abierta.

%-------------------------------------------------------------------------
% \section*{Reproducibilidad}
% Opcional: enlace al notebook / repositorio.

\end{document}
